% Module 6: Quantum Communication
% Estimated slides: 30 | Duration: ~2 hours
% Key topics: Quantum teleportation, superdense coding, BB84, QKD,
%             quantum repeaters, quantum internet overview
% Labs: Teleportation circuit in Qiskit; BB84 simulation

%%%%%%%%%%%%%%%%%%%%%%%%%%%%%%%%%%%%%%%%%%%%%%%%%%%%%%%%%%
\begin{frame}[fragile]\frametitle{}
\begin{center}
{\Large Module 6: Quantum Communication}
\end{center}
\end{frame}

%%%%%%%%%%%%%%%%%%%%%%%%%%%%%%%%%%%%%%%%%%%%%%%%%%%%%%%%%%%
\begin{frame}[fragile]\frametitle{Quantum vs Classical Communication}
\begin{itemize}
\item Classical communication: send bits over a wire, fiber, or radio wave
\item Quantum communication: transmit quantum states (qubits) between parties
\item Key quantum advantage: security guaranteed by physics, not computational hardness
\item Eavesdropping on quantum channels is detectable due to measurement collapse
\item Applications: quantum key distribution, quantum networks, distributed quantum computing
\item Quantum internet: a network that transmits quantum states between nodes
\end{itemize}
% Suggested diagram: Classical network vs quantum network with EPR pair links
\end{frame}

%%%%%%%%%%%%%%%%%%%%%%%%%%%%%%%%%%%%%%%%%%%%%%%%%%%%%%%%%%%
\begin{frame}[fragile]\frametitle{Quantum Teleportation: What It Is}
\begin{itemize}
\item Quantum teleportation transfers an unknown qubit state from Alice to Bob
\item Does NOT transmit matter or energy faster than light
\item Requires: a shared entangled pair (Bell state) + a classical communication channel
\item The original qubit is destroyed in the process (no-cloning is respected)
\item Why ``teleportation''? The state ``appears'' at Bob without traveling through space
\item Not useful for communication; useful for distributed quantum computing
\end{itemize}
% Suggested diagram: Alice-Bob teleportation setup with EPR pair and classical channel
\end{frame}

%%%%%%%%%%%%%%%%%%%%%%%%%%%%%%%%%%%%%%%%%%%%%%%%%%%%%%%%%%%
\begin{frame}[fragile]\frametitle{Quantum Teleportation: Step-by-Step Protocol}
\begin{itemize}
\item \textbf{Setup}: Alice has qubit $|\psi\rangle$ to send; Alice and Bob share Bell state $|\Phi^+\rangle$
\item \textbf{Step 1}: Alice applies CNOT between $|\psi\rangle$ and her half of the Bell pair
\item \textbf{Step 2}: Alice applies H to $|\psi\rangle$
\item \textbf{Step 3}: Alice measures her two qubits; gets one of 4 outcomes (00, 01, 10, 11)
\item \textbf{Step 4}: Alice sends the 2 classical bits to Bob
\item \textbf{Step 5}: Bob applies corrections (X, Z) based on Alice's bits
\item \textbf{Result}: Bob's qubit is now in state $|\psi\rangle$; Alice's original is destroyed
\end{itemize}
% Suggested diagram: Full teleportation circuit diagram (3 qubits: psi, Alice's EPR half, Bob's EPR half)
\end{frame}

%%%%%%%%%%%%%%%%%%%%%%%%%%%%%%%%%%%%%%%%%%%%%%%%%%%%%%%%%%%
\begin{frame}[fragile]\frametitle{Why Teleportation Doesn't Violate Relativity}
\begin{itemize}
\item Alice must send 2 classical bits to Bob before he can recover the state
\item Classical bits travel at most at the speed of light
\item Without the classical bits, Bob's qubit is just maximally mixed noise
\item The quantum state doesn't travel faster than light --- only after classical bits arrive
\item No information is transmitted by the entanglement alone
\item Teleportation is still remarkable: the state is perfectly transferred with 0 bits of overhead about the state itself
\end{itemize}
\end{frame}

%%%%%%%%%%%%%%%%%%%%%%%%%%%%%%%%%%%%%%%%%%%%%%%%%%%%%%%%%%%
\begin{frame}[fragile]\frametitle{Lab: Teleportation Circuit in Qiskit}
\begin{itemize}
\item Implement the full 3-qubit teleportation protocol
\item Prepare an arbitrary state $|\psi\rangle$ on qubit 0 (use $R_y(\pi/3)$)
\item Share Bell state on qubits 1 and 2
\item Execute protocol; verify Bob's qubit matches the original state
\end{itemize}
\begin{verbatim}
qc = QuantumCircuit(3, 2)
qc.ry(0.9, 0)          # prepare state to teleport
qc.h(1); qc.cx(1, 2)  # Bell pair
qc.cx(0, 1); qc.h(0)  # Alice's operations
qc.measure([0,1], [0,1])
# Classical corrections on qubit 2 based on measurements
\end{verbatim}
% Suggested diagram: Teleportation circuit with labeled stages
\end{frame}

%%%%%%%%%%%%%%%%%%%%%%%%%%%%%%%%%%%%%%%%%%%%%%%%%%%%%%%%%%%
\begin{frame}[fragile]\frametitle{Superdense Coding}
\begin{itemize}
\item Superdense coding: send 2 classical bits using just 1 qubit (+ pre-shared entanglement)
\item Alice and Bob share a Bell state $|\Phi^+\rangle$ in advance
\item Alice applies one of 4 local operations (I, X, Z, XZ) to encode 2 classical bits
\item Alice sends her qubit to Bob
\item Bob performs a Bell measurement and decodes 2 bits
\item Achieves capacity of 2 bits per qubit (Holevo bound saturated with entanglement)
\end{itemize}
% Suggested diagram: Superdense coding circuit with Alice's 4 operations for 00, 01, 10, 11
\end{frame}

%%%%%%%%%%%%%%%%%%%%%%%%%%%%%%%%%%%%%%%%%%%%%%%%%%%%%%%%%%%
\begin{frame}[fragile]\frametitle{Quantum Key Distribution: The Security Problem}
\begin{itemize}
\item Classical cryptography: security based on computational hardness (RSA, ECC)
\item Threat: Shor's algorithm on a large quantum computer can break RSA and ECC
\item Quantum Key Distribution (QKD): security based on laws of physics, not computation
\item Any eavesdropping on a quantum channel disturbs the signal --- detectable
\item QKD allows Alice and Bob to detect eavesdroppers and abort if compromised
\item Only provides key exchange, not encryption itself
\end{itemize}
% Suggested diagram: Alice-Bob-Eve triangle with quantum channel and classical channel
\end{frame}

%%%%%%%%%%%%%%%%%%%%%%%%%%%%%%%%%%%%%%%%%%%%%%%%%%%%%%%%%%%
\begin{frame}[fragile]\frametitle{BB84 Protocol: Step by Step}
\begin{itemize}
\item Alice generates random bits and random bases (Z or X basis) for each bit
\item Alice encodes each bit as a qubit in her chosen basis and sends to Bob
\item Bob chooses a random basis for each qubit and measures
\item Bob and Alice compare bases over a classical channel (public)
\item They keep only bits where they chose the same basis (~50\%)
\item If Eve measured qubits (random basis), she introduces errors ($\sim$25\% error rate)
\item Alice and Bob check a sample of bits: if error rate too high, abort
\end{itemize}
% Suggested diagram: BB84 table showing Alice's bits/bases, Bob's bases/results, sifting
\end{frame}

%%%%%%%%%%%%%%%%%%%%%%%%%%%%%%%%%%%%%%%%%%%%%%%%%%%%%%%%%%%
\begin{frame}[fragile]\frametitle{BB84: Why Eavesdropping is Detectable}
\begin{itemize}
\item Eve must guess which basis to use for each qubit (50\% chance of correct guess)
\item Wrong basis choice: Eve collapses the qubit, then re-sends a random state
\item When Alice and Bob compare: 25\% of their matching-basis bits will disagree
\item Error threshold: $< 11\%$ QBER (Quantum Bit Error Rate) considered secure
\item Above threshold: channel is compromised, discard key and retry
\item Security proof does not depend on Eve's computational power
\end{itemize}
% Suggested diagram: QBER illustration with and without eavesdropper
\end{frame}

%%%%%%%%%%%%%%%%%%%%%%%%%%%%%%%%%%%%%%%%%%%%%%%%%%%%%%%%%%%
\begin{frame}[fragile]\frametitle{Lab: BB84 Simulation}
\begin{itemize}
\item Simulate the full BB84 protocol in Python (no quantum hardware needed)
\item Generate random bits and bases for Alice and Bob
\item Simulate Eve intercept-and-resend attack
\item Calculate QBER with and without Eve
\item Observe: Eve's presence significantly increases error rate
\end{itemize}
\begin{verbatim}
import numpy as np
alice_bits = np.random.randint(2, size=100)
alice_bases = np.random.randint(2, size=100)  # 0=Z, 1=X
# Encode, transmit, measure, sift...
\end{verbatim}
\end{frame}

%%%%%%%%%%%%%%%%%%%%%%%%%%%%%%%%%%%%%%%%%%%%%%%%%%%%%%%%%%%
\begin{frame}[fragile]\frametitle{Real QKD Systems Today}
\begin{itemize}
\item Commercial QKD products: ID Quantique (Switzerland), Toshiba, QuantumCTek (China)
\item Distance limits: fiber QKD ~100-600 km; free-space to satellites > 1000 km
\item China's Micius satellite: QKD over 1200 km (2017)
\item China's Quantum Internet: 2000 km backbone network (2021)
\item Limitation: cannot amplify quantum signals (no-cloning) --- need quantum repeaters
\item Quantum repeaters: use entanglement swapping to extend range (active research)
\end{itemize}
% Suggested diagram: Global map of operational QKD networks
\end{frame}

%%%%%%%%%%%%%%%%%%%%%%%%%%%%%%%%%%%%%%%%%%%%%%%%%%%%%%%%%%%
\begin{frame}[fragile]\frametitle{The Quantum Internet Vision}
\begin{itemize}
\item A global network that can transmit quantum states (not just classical bits)
\item Enables: perfectly secure communication, distributed quantum computing, quantum sensing
\item Stages (Wehner et al., Science 2018): trusted repeater, entanglement distribution, quantum memory, fault-tolerant
\item Current reality: Stage 1-2 (fiber QKD networks, satellite QKD demonstrated)
\item Quantum repeaters require quantum memory (still in development)
\item Timeline to full quantum internet: 10-20 years
\end{itemize}
% Suggested diagram: Quantum internet development stages diagram
\end{frame}

%%%%%%%%%%%%%%%%%%%%%%%%%%%%%%%%%%%%%%%%%%%%%%%%%%%%%%%%%%%
\begin{frame}[fragile]\frametitle{Post-Quantum Cryptography (PQC)}
\begin{itemize}
\item QKD secures future communication, but what about data encrypted today?
\item ``Harvest now, decrypt later'': adversaries record encrypted data now, decrypt with quantum computer later
\item NIST PQC Standardization: algorithms resistant to quantum attacks using classical computers
\item Standards finalized 2024: CRYSTALS-Kyber (key exchange), CRYSTALS-Dilithium (signatures), SPHINCS+
\item Based on hard problems: lattice problems, hash functions, not RSA/ECC
\item Migration timeline: all systems should upgrade to PQC by 2030 (NIST recommendation)
\end{itemize}
% Suggested diagram: PQC algorithm families and their quantum security claims
\end{frame}

%%%%%%%%%%%%%%%%%%%%%%%%%%%%%%%%%%%%%%%%%%%%%%%%%%%%%%%%%%%
\begin{frame}[fragile]\frametitle{Module 6 Summary}
\begin{itemize}
\item Quantum teleportation: transfers qubit state using entanglement + 2 classical bits
\item Superdense coding: sends 2 classical bits using 1 qubit + pre-shared entanglement
\item Neither violates relativity: classical channel is required
\item BB84: first QKD protocol; security based on physics, not computation
\item Eavesdropping is detectable via quantum bit error rate (QBER)
\item Real QKD systems deployed today; quantum internet is being built
\item Post-quantum cryptography: classical algorithms safe against quantum attacks
\end{itemize}
\end{frame}
